\section{Implementation}
\label{sec:impl}

This section walks through the pipeline stage by stage. Code excerpts are drawn
from the extension's source and lightly trimmed for space.

\subsection{Extension Structure and Activation}
The extension is a Manifest~V3 package with no content scripts. Its content
security policy permits \code{wasm-unsafe-eval} so the bundled Tesseract
WebAssembly core can run, and the OCR engine, language data, PDF.js and pdf-lib
are all bundled---no CDN fetch occurs at runtime.

Activation is on demand (decision~\ref{dec:ondemand}). Two entry points---the
toolbar icon (\code{action.onClicked}) and a keyboard shortcut
(\code{commands.onCommand}, \code{Ctrl+Shift+F} on every platform)---run one
handler (Listing~\ref{lst:activate}). On a searchable copy it restores the
original, whose URL is carried in the copy's \code{src} parameter. Otherwise it
starts the offscreen document and sends it the job, including the copy URL and
cache key the service worker has chosen, so that the naming of copies lives in
one place.

\begin{lstlisting}[caption={On-demand activation: one handler for the toolbar icon and the shortcut. It restores the original when the tab shows a copy, and otherwise starts a job (\code{service-worker.js}).}, label={lst:activate}]
async function toggle(tab) {
  // Showing a searchable copy? Restore the original PDF.
  if (isCopyUrl(tab.url)) {
    const original = originalOf(tab.url);
    await dropSwap(tab.id);
    await deleteCopy(tab.url);
    if (original) chrome.tabs.update(tab.id, { url: original });
    return;
  }
  if (!isFetchableUrl(tab.url)) return badgeError(tab.id, "open a PDF in this tab first");

  badge(tab.id, "...", "Making this PDF searchable...");
  await ensureOffscreen();
  const url = tab.url.split("#")[0];
  const copyUrl = newCopyUrl(url);   // chrome-extension://<id>/searchable/<uuid>/<name>?src=<url>
  chrome.runtime.sendMessage({ type: "fii-process", tabId: tab.id, url,
    lang: await getLang(), copyUrl, cache: COPY_CACHE, cacheKey: copyKey(copyUrl) });
}
chrome.action.onClicked.addListener((tab) => toggle(tab));
chrome.commands.onCommand.addListener((cmd, tab) => { if (cmd === "toggle-search") toggle(tab); });
\end{lstlisting}

Local \code{file://} PDFs additionally require the per-extension ``Allow access to
file URLs'' setting; if it is off, the handler shows a red badge explaining this
and opens the extension's details page, where the toggle lives.

\subsection{Download, Validation, and Identity}
The offscreen document downloads the PDF again from the tab's URL with
\code{credentials: "include"}, so that PDFs on sites the user is signed in to
return the same file the tab shows. The trigger accepts \emph{any} web or file
URL, because many PDFs are served from addresses without ``.pdf'' (for example
\code{arxiv.org/pdf/2111.03017} or \code{download.php?id=7}). Validation therefore
happens on the bytes: a PDF must contain the signature \code{\%PDF-} within its
first kilobyte. A page that returns HTML is reported as ``not a PDF (or a
sign-in page)''.

The bytes are then loaded twice, by PDF.js (for analysis and rendering) and by
pdf-lib (for writing), while \code{crypto.subtle.digest} computes their SHA-256 in
parallel. If either library reports encryption, the job ends with a clear
message: pdf-lib cannot write into encrypted files. pdf-lib's
\code{EncryptedPDFError} is transpiled ES5, so \code{instanceof} does not
recognize it and the check matches its message instead. The PDF.js loading task
is destroyed in a \code{finally} block, because the offscreen document lives all
session and would otherwise keep every processed document in memory.

\subsection{Text Detection}
For each page the processor first calls \code{getTextContent()}. OCR runs when
the page paints raster images \emph{or} when no text is extractable at all (a page
may be a scan, or have its text converted to vector outlines---both invisible to
\code{getTextContent()}). Image detection inspects the operator list:

\begin{lstlisting}[caption={Image detection: scans the PDF.js operator list for raster paint opcodes to decide whether a page needs OCR (\code{processor.js}).}, label={lst:hasimages}]
const IMAGE_OPS = new Set([
  pdfjsLib.OPS.paintImageXObject,      pdfjsLib.OPS.paintInlineImageXObject,
  pdfjsLib.OPS.paintImageMaskXObject,  pdfjsLib.OPS.paintImageXObjectRepeat,
]);
async function pageHasImages(page) {
  const opList = await page.getOperatorList();
  return opList.fnArray.some((fn) => IMAGE_OPS.has(fn));
}
\end{lstlisting}

If image detection fails it \emph{fails open}---assuming images are present---on
the principle that an unnecessary OCR pass is cheaper than silently missing text.

\subsection{OCR Rendering and Preprocessing}
Pages that need OCR are rendered to an offscreen canvas about 3000\,px wide
($\approx$360\,DPI for a Letter page), with PDF.js's \code{print} intent so that
rendering is not paced by \code{requestAnimationFrame}, which never fires in a
hidden document. The canvas is released immediately after recognition.

Before recognition, the image is converted to grayscale and preprocessed to
attack the two hardest cases, low contrast and blur:
\begin{enumerate}
  \item \textbf{Gain-capped contrast stretch.} With $P_{2}$ and $P_{98}$ the 2nd
        and 98th percentiles of the luminance histogram, every gray level $v$ is
        remapped as
        \begin{equation}
          v' = 255 - 255\,\frac{P_{98} - v}{R},
          \qquad
          R = \max\bigl(R_{\min},\; P_{98} - P_{2}\bigr),
          \quad R_{\min} = 96 .
          \label{eq:stretch}
        \end{equation}
        When the percentiles are far apart this is the classic stretch
        ($P_2\mapsto 0$, $P_{98}\mapsto 255$). The floor $R_{\min}$ matters on
        mostly blank pages: when text covers less than 2\% of the pixels, $P_2$
        is itself background. On the first test pages $P_2$ was 233--238 against
        $P_{98}=240$, so the uncapped stretch multiplied contrast by up to
        $127\times$, turned faint JPEG noise into black speckle, and made one page
        unreadable. Anchoring at $P_{98}$ keeps the background white, and the
        floor limits the gain to $255/96 \approx 2.7$.
  \item \textbf{Unsharp mask.} Edges are re-sharpened as
        $\text{sharp} = g + \alpha\,(g - g_{\text{blur}})$ with $\alpha=1$, where
        $g_{\text{blur}}$ comes from a separable box blur implemented with a
        sliding-window sum, so its cost is $O(n)$ independent of radius.
\end{enumerate}
Tesseract performs its own binarization afterward, so this stage only needs to
hand it a clean, high-contrast grayscale image.

\subsection{OCR Engine Orchestration}
Recognition runs through a Tesseract.js \emph{scheduler} that distributes pages
across a pool of Web Workers. The worker count scales with CPU cores but is
capped, because each worker loads its own copy of the language model:
\begin{equation}
  N_{\text{workers}} = \max\!\big(1,\ \min(4,\ \texttt{hardwareConcurrency}-1)\big).
\end{equation}
The processor keeps $N_{\text{workers}}$ pages in flight at once. JavaScript is
single-threaded, so pdf-lib edits to different pages interleave safely; only the
rendering and recognition awaits overlap. Workers use \emph{sparse text} page
segmentation, which finds table cells, captions, and scattered text far better
than the default layout analysis. Three language modes are supported---\code{eng},
\code{dan}, and \code{eng+dan}---with one worker pool per language. Sixty seconds
after the job queue empties, all pools are terminated to free their memory; they
are recreated on the next job.

\subsection{Post-Processing for Findability}
Because native find performs exact substring matching, fuzzy matching is
impossible at search time; the post-processing stage improves the \emph{stored}
text instead. It drops near-zero-confidence noise, trims non-letter junk from
word edges (with Unicode letter classes, so Danish words keep their
\textit{æ}, \textit{ø} and \textit{å}), and emits a \emph{correction variant}
for words that look like digit/letter confusions (\code{he11o}\,$\to$\,\code{hello},
\code{ca5h}\,$\to$\,\code{cash}). Tokens that are more than half digits are left
alone as genuine numbers or codes:

\begin{lstlisting}[caption={Emitting a correction variant for OCR tokens with likely digit/letter confusions (\code{ocr-postprocess.js}).}, label={lst:variant}]
const DIGIT_TO_LETTER = { "0":"o","1":"l","5":"s","6":"b","8":"b","9":"g","2":"z" };
out.push({ text, bbox: w.bbox });                 // the word as recognized
const variant = correctionVariant(text);          // null unless mostly letters + some digits
if (variant && variant !== text)
  out.push({ text: variant, bbox: w.bbox, variant: true });
\end{lstlisting}

The earlier DOM viewer injected both forms and hid the variant from selection.
Inside a real PDF there is no such distinction: two words at one spot would
double Chrome's match count and duplicate copied text. The processor therefore
embeds a single form, preferring the corrected variant, since it is what the
image most likely says.

\subsection{Writing the Invisible Layer}
Each recognized word is mapped into PDF user space as described in
Section~\ref{sec:coords} and de-duplicated against the page's existing text.
The survivors are appended to the page's content stream as one invisible text
object each (Listing~\ref{lst:embed}). The whole block is wrapped in
\code{q}/\code{Q} so its text state cannot leak into anything drawn later.
Helvetica's WinAnsi encoding covers English and Danish; a word containing a
character it cannot encode is skipped rather than failing the page.

\begin{lstlisting}[caption={Writing one invisible, width-matched text object per word with pdf-lib (\code{embed.js}).}, label={lst:embed}]
for (const wd of words) {
  const encoded = font.encodeText(wd.text);
  const natural = font.widthOfTextAtSize(wd.text, wd.h);
  const squeeze = Math.max(1, Math.min(1000, (wd.w / natural) * 100));  // Tz, percent
  const cos = Math.cos(wd.angle), sin = Math.sin(wd.angle);
  ops.push(
    PDFLib.beginText(),
    PDFLib.setTextRenderingMode(PDFLib.TextRenderingMode.Invisible),     // 3 Tr
    PDFLib.setFontAndSize(fontKey, wd.h),
    PDFLib.setCharacterSqueeze(squeeze),
    PDFLib.setTextMatrix(cos, sin, -sin, cos, wd.x, wd.y),
    PDFLib.showText(encoded),
    PDFLib.endText(),
  );
}
page.pushOperators(PDFLib.pushGraphicsState(), ...ops, PDFLib.popGraphicsState());
\end{lstlisting}

The document is then saved with its metadata untouched, and the bytes are
stored in Cache Storage as an \code{application/pdf} response under the key the
service worker supplied.

\subsection{Caching}
Recognized words are normalized to the unit square (Eq.~\eqref{eq:normalize}) and
stored in IndexedDB under \code{<sha256>::<lang>::<page>}. Because the key is the
file's content hash, the same PDF reached through two different URLs is
recognized once, and a different file at a familiar URL is never matched to
stale words. Empty results are cached too, so text-only pages are not re-examined.
Each record carries a \emph{format version}, bumped whenever the OCR pipeline
changes, and a last-used timestamp. A cache hit refreshes the timestamp, and
after every job the oldest records are deleted until at most 2{,}000 pages
remain.

\subsection{Serving and Cleaning Up the Copy}
When the job is done, the service worker checks that the tab still shows the
same PDF (the user may have moved on), records which copy the tab shows, and
navigates it to the copy URL. Chrome requests that URL from the extension, and
the service worker's \code{fetch} handler answers it (Listing~\ref{lst:serve}).
Chrome sees an \code{application/pdf} response and opens its own viewer; because
the path ends in the original file name, the tab title is unchanged. The success
badge is set only after the tab finishes loading, because Chrome clears a tab's
badge whenever it navigates.

\begin{lstlisting}[caption={Serving searchable copies at the extension's own URLs. Cache Storage only accepts http(s) keys, so copies are stored under a placeholder origin with the same path (\code{service-worker.js}).}, label={lst:serve}]
self.addEventListener("fetch", (event) => {
  const url = new URL(event.request.url);
  if (url.origin !== self.location.origin || !url.pathname.startsWith("/searchable/")) return;
  event.respondWith(serveCopy(url));
});
async function serveCopy(url) {
  const hit = await (await caches.open(COPY_CACHE)).match(COPY_KEY_ORIGIN + url.pathname);
  if (hit) return hit;
  // The copy is gone (e.g. the browser restarted): fall back to the original PDF.
  const src = url.searchParams.get("src");
  if (src && isFetchableUrl(src)) return Response.redirect(src, 302);
  return new Response("This searchable copy is no longer available.", { status: 404 });
}
\end{lstlisting}

Copies are deleted when the user restores the original, when their tab
navigates elsewhere or closes, and when the browser starts. A tab that returns to
a deleted copy---through session restore or the browser's history---simply
redirects to the original PDF.
